\section*{Limitations}
\label{sec:limitations}

\paragraph{Corpus, scale and generalization.}
The evaluation covers one English public-domain series, five volumes and one generator.
The series informed development, and all five volumes fit within the generator's context
window. Cutoffs share characters and evidence, so their estimates do not represent
independent works. Results may differ for other genres, languages, models or series
exceeding one context window. Commercial development series contribute no reported
measurements. A held-out series and repeated generations are needed to assess
generalization and variability.

\paragraph{Annotation and gold evidence.}
Gold facts and evaluation labels were finalized by Claude, who knew system identity and
also assisted in implementation. A common finalizer does not establish independent or
blinded evaluation. The pair-context verdicts come from Claude; alternative provisional
pair judgments yield precision \res{REV-M3-pair2} rather than
\respt{M3-anne-avg-B5}. Thus judgment sensitivity can affect comparisons across systems and
text conditions even with one finalizing annotator.
\ifnwhuman In the earlier review, one human accepted renewed LLM answers for all gold facts and 10\%
samples of recall, precision and citation labels, with those answers visible. This provides
no independent agreement estimate for that pass. The later check recomputes agreement between
supplied human final-label sets and with an accepted AI-assisted reference (Tables~\ref{tab:human-agreement}
and~\ref{tab:human-ai-agreement}), and joint nominal $\alpha$ across those three label sets
(Table~\ref{tab:human-alpha}). Both human label sets were produced entirely by human
reviewers working from scratch, without the AI answers (Appendix~\ref{app:annotation}).
Agreement does not establish annotation accuracy. Both checks exclude disclosure and
treated-text judgments.\else The planned human review is not part of this version.\fi{}
The audit also found that some gold evidence paragraphs establish their facts only partially
or not at all (Section~\ref{sec:setup}). The inventory is retained; resolving these cases
requires reviewing the underlying annotation data. Precision and citation estimates use
40 sampled atoms per cell and remain sensitive to sampling and judgment rules.

\paragraph{Incomplete disclosure assessment.}
M4 examines a selected inventory of later facts through candidate retrieval, rather than
exhaustively reading every assertion for every possible spoiler. Validated seeded fixtures
give approximately 85\% retrieval for named paraphrases and 70\% for named implications,
with substantially lower recall for pronouns on other pages. The fixtures were written
and validated by the same LLM annotator. Missed disclosures, unassessed plot points,
identity changes and surfaces outside the published directory can therefore escape the
evaluation. Different artifact sizes and clustered assertions further limit the
interpretation of zero counts: a zero or low count does not certify spoiler-free output.

\paragraph{Development feedback and single runs.}
The reported \sys is a second build revised after first-build labels became available;
baselines were not retuned in that round. Table~\ref{tab:builds} reports both builds.
The complete prefix-versus-full graph comparison used first-build code. For the reported
build, prefix-discovered indexes and wikis were evaluated on the original text only,
with one-pass LLM labels (Appendix~\ref{app:additional}). Entity-index differences include
candidate-mining rules and stochastic model decisions. Each system and condition was
generated once; reported standard deviations reflect cutoff variation, not run-to-run
or model variation. Some original-text citation labels were added in later sessions.
B4 differs from \sys in page scope and generation as well as citation attachment, so
that comparison does not isolate their effects, and the repairs suggested in
Section~\ref{sec:discussion} are untested.
\ifdefined\nwvtwo The gains of Section~\ref{sec:results-gains} are measured against
comparison pages written in one call under a sentence budget and, for the retrieval
systems, a 20,000-token passage budget; longer or multi-call pages were not tried. The
split into central and other characters was made after the results existed, and the count
of facts lost across cutoffs also counts labeling inconsistencies between adjacent
cutoffs.\fi

\paragraph{Recognition intervention.}
High closed-book recall shows that the generator recognizes this public-domain work;
it does not establish how much each text-based system relies on memory.
\ifnwdecon Renaming and paraphrasing sharply reduce closed-book recall, but neither
remove recognition completely nor preserve the text perfectly. The treatment was not
proofread against its source; \res{DECON-residual-forms} paragraphs retain inflected
source names, and name collisions affect page selection. Some outputs recover the
island's original name despite its absence from the treated text. Gold fidelity was
checked by an unblinded LLM, without repeating the stricter paragraph-sufficiency audit.
The treatment changes wording, entity resolution and potentially meaning, limiting
cross-condition comparisons even with labels finalized by the same annotator.
\fi{} A work published after the generator's training cutoff would offer a complementary
test of recognition, as in NoCha \citep{karpinska2024nocha}.

\paragraph{Temporal and discourse representation.}
Cutoffs are finished volumes. Two conflicting attribute values within one volume are
treated as an extraction conflict, although both may describe genuine changes. Identity
links are not dated, and predictions or beliefs can become factual relations if attribution
or modality is lost. Attribute intervals track explicit replacements rather than implicit
aging: an early age can remain displayed after later events (Appendix~\ref{app:qualitative}).
The system also lacks systematic cross-claim inference and may miss pronoun links across
extraction windows. These qualitative observations have no measured prevalence
(Appendix~\ref{app:observations}).

\paragraph{Reader utility and publication scope.}
We measure factual coverage, support, citation sufficiency and disclosure, but not readers'
ability to find information or their experience of spoilers. The additional linked pages
may provide value or impose navigation costs; no reader study establishes either.
The system omits portraits because development trials could not reliably ground character
illustrations in textual evidence. External search engines and deliberately opened later
cutoffs are outside the audited artifact.
